%!TEX root=../../note.tex
\subsection{Subsets and Set Equality}
To compare two sets we ask whether every element of one belongs to the other.
This gives the subset relation, and two subset relations together give set
equality.

\subsubsection{Disjoint Sets and Subsets}
\begin{definition}[Disjoint]
    Two sets $S$ and $T$ are \emph{disjoint} when they have no element in
    common, that is,
    \begin{equation*}
        S \cap T = \emptyset.
    \end{equation*}
\end{definition}

\begin{definition}[Subset]
    A set $S$ is a \emph{subset} of a set $T$, written $S \subseteq T$, when
    every element of $S$ belongs to $T$.
\end{definition}

\begin{definition}[Proper Subset]
    A set $S$ is a \emph{proper subset} of a set $T$, written
    $S \subsetneq T$, when $S \subseteq T$ and some element of $T$ does not
    belong to $S$.
\end{definition}

\begin{remark}
    If $S \subseteq T$, we also say $T$ is a \emph{superset} of $S$. If
    $S \subsetneq T$, we say $T$ is a \emph{proper superset} of $S$. We
    write $S \nsubseteq T$ when $S$ is not a subset of $T$, that is, when
    some element of $S$ does not belong to $T$.
\end{remark}

\begin{example}
    \begin{itemize}
        \item $S \subseteq S$ for every set $S$.
        \item $\emptyset \subseteq S$ for every set $S$: the empty set has no
        element that could fail to belong to $S$.
        \item $\N \subseteq \Z$; in fact $\N \subsetneq \Z$, since $0\in\Z$
        but $0\notin\N$.
        \item $\Z \nsubseteq \N$, since $-1\in\Z$ but $-1\notin\N$.
    \end{itemize}
\end{example}
% TODO(check): the example list ended with an empty item in lecture notes; one example may be missing.

\subsubsection{Proving a Subset Relation}
To prove that $S \subseteq T$, prove the universally quantified implication
\begin{equation*}
    \forall x \in \mathcal{U},\ \paren{x \in S} \implies \paren{x \in T},
\end{equation*}
where $\mathcal{U}$ is the universe of discourse: let $x$ be an arbitrary
element of $S$ and show that $x\in T$. To prove $S \nsubseteq T$, one
counterexample suffices: an element of $S$ that is not in $T$.

\begin{example}
    Let $A = \set{n \in \N : 4 \mid (n - 3)}$ and
    $B = \set{2k + 1 : k \in \Z}$.
    \begin{enumerate}[label=(\roman*)]
        \item Prove that $A \subseteq B$.
        \item Prove that $A \subsetneq B$.
    \end{enumerate}
\end{example}
\begin{remark}[Proof idea]
    Elements of $B$ are the odd integers. For (i), unpack $4\mid(n-3)$ into
    $n=4k+3$ and rewrite it in the form $2(\cdot)+1$. For (ii), find an odd
    integer that is not in $A$.
\end{remark}
\begin{proof}
    (i) Let $n \in \N$ and assume $n \in A$. By the definition of $A$,
    $4 \mid \paren{n - 3}$, so there is some $k \in \Z$ such that
    $n - 3 = 4k$. Then
    \begin{equation*}
        n = 4k + 3 = 4k + 2 + 1 = 2\paren{2k + 1} + 1.
    \end{equation*}
    Since $2k + 1 \in \Z$, we have $n \in B$. Therefore $A \subseteq B$.

    (ii) By (i), $A\subseteq B$. Now $1 = 2\cdot 0 + 1 \in B$, but
    $1 - 3 = -2$ and $4 \nmid -2$, so $1 \notin A$. Hence $B$ has an element
    that is not in $A$, and $A \subsetneq B$.
\end{proof}

\subsubsection{Proving Set Equality}
Two sets are equal when they have exactly the same elements. So, to prove
$S = T$, we prove both
\begin{equation*}
    S \subseteq T \qquad\text{and}\qquad T \subseteq S.
\end{equation*}

\begin{example}
    Let $S = \set{-1, 0, 1}$ and $T = \set{x \in \R : x^3 = x}$. Prove that
    $S = T$.
\end{example}
\begin{proof}
    \textbf{$S \subseteq T$.} Let $x \in S$. Then $x$ is $-1$, $0$, or $1$,
    and we check each case:
    \begin{equation*}
        \paren{-1}^3 = -1, \qquad 0^3 = 0, \qquad 1^3 = 1.
    \end{equation*}
    In each case $x^3 = x$, so $x \in T$.

    \textbf{$T \subseteq S$.} Let $x \in T$. Then $x^3 = x$, so
    \begin{equation*}
        0 = x^3 - x = x\paren{x^2 - 1} = x\paren{x - 1}\paren{x + 1}.
    \end{equation*}
    A product of real numbers is $0$ only if one factor is $0$, so $x = 0$,
    $x = 1$, or $x = -1$. Thus $x \in S$.

    Therefore $S = T$.
\end{proof}

\begin{example}
    Prove that for all sets $S$ and $T$, $S = T$ if and only if
    $S \cap T = S \cup T$.
\end{example}
\begin{remark}[Proof idea]
    Each direction is a set equality, proved by two subset relations. For
    $(\Leftarrow)$, an element of $S$ lies in $S \cup T$, hence in
    $S \cap T$, hence in $T$; the same chain works with $S$ and $T$
    swapped.
\end{remark}
\begin{proof}
    $(\Rightarrow)$ Suppose $S = T$. We show $S \cap T = S \cup T$.

    First, $S \cap T \subseteq S \cup T$: if $x \in S \cap T$, then
    $x \in S$ and $x \in T$; in particular $x \in S$, so $x \in S \cup T$.

    Next, $S \cup T \subseteq S \cap T$: if $x \in S \cup T$, then
    $x \in S$ or $x \in T$. Since $S = T$, in either case $x \in S$ and
    $x \in T$, so $x \in S \cap T$.

    (Alternatively, substitute $T=S$: $S \cap T = S \cap S = S = S \cup S = S \cup T$.)

    $(\Leftarrow)$ Suppose $S \cap T = S \cup T$. We show $S = T$.

    First, $S \subseteq T$: for $x\in S$,
    \begin{equation*}
        x \in S \implies x \in S \cup T \implies x \in S \cap T \implies x \in T.
    \end{equation*}

    Second, $T \subseteq S$: for $x\in T$,
    \begin{equation*}
        x \in T \implies x \in S \cup T \implies x \in S \cap T \implies x \in S.
    \end{equation*}
    The middle step in each chain uses the hypothesis $S \cup T = S \cap T$.

    Therefore $S = T$.
\end{proof}

\subsubsection*{In practice}
\begin{itemize}
    \item \textbf{Element chasing.} To prove $S\subseteq T$, write ``Let
    $x\in S$'', unpack the definition of $S$, and rebuild the definition of
    $T$. To disprove it, exhibit one element of $S$ not in $T$.
    \item \textbf{Set equality.} Prove $S\subseteq T$ and $T\subseteq S$
    separately. For a finite set such as $\set{-1,0,1}$, the first inclusion
    is a finite case check.
    \item \textbf{Proper subset.} Prove $S\subseteq T$, then exhibit one
    element of $T$ outside $S$.
\end{itemize}
